\begin{abstract}
云服务商迫切需要高效地服务大语言模型（LLM）；在处理每个请求后缓存其中间结果（{\kvcache}），能够显著改善服务吞吐量与延迟。然而，人们对于 LLM 服务究竟如何受益于 {\kvcache} 缓存仍缺乏充分认识，而缓存淘汰策略等系统设计决策又高度依赖工作负载。

本文首次对一家领先 LLM 服务商的 {\kvcache} 工作负载模式进行了系统刻画。相较于聚焦合成工作负载的既有研究，我们得到了一些此前未被覆盖的观察：{\kvcache} 的复用在请求之间呈偏斜分布，单轮请求之间的复用与多轮请求同样重要；从所有请求整体看，复用时间与复用概率具有多样性，但对某一具体请求类别而言，其模式往往可预测；为达到理想缓存命中率所需的总体缓存容量较为适中。基于这些特征，我们进一步提出一种工作负载感知的缓存淘汰策略；在真实轨迹下，尤其当缓存容量受限时，该策略能够改善服务性能。
\end{abstract}
